% Conjectura manuscript, CJ-0001 (upper half of the Collins--Shor conjecture).
% DRAFT, NOT REVIEWED BY A HUMAN, NOT FOR DISTRIBUTION.
% Prepared on 2026-11-27 by the publication agent (Claude, model
% claude-opus-5-5) from the following project files only:
%   research/active/CJ-0001-collins-shor-peak/notes/upper-bound.md (revision 1,
%     commit 2037b18, SHA-256 ad21eeb615ab3a50...49e0ed)
%   research/active/CJ-0001-collins-shor-peak/notes/path-peak.md (commit 2037b18,
%     SHA-256 b65a9931ed849704...e78016)
%   research/active/CJ-0001-collins-shor-peak/reviews/2026-11-27-*upper-bound*.md,
%     2026-11-27-*path-peak*.md
%   research/candidates/CJ-0001-collins-shor-peak.md (literature audits)
% The mathematical content is intended to be identical to the notes; defects
% listed in the reviews are folded in as wording, citation or written-out steps.
% Plain LaTeX with standard packages only.
\documentclass[11pt]{amsart}
\usepackage{amsmath,amssymb,amsthm}
\PassOptionsToPackage{hyphens}{url}
\usepackage[colorlinks=true,allcolors=blue]{hyperref}

\newtheorem{theorem}{Theorem}[section]
\newtheorem{lemma}[theorem]{Lemma}
\newtheorem{proposition}[theorem]{Proposition}
\newtheorem{corollary}[theorem]{Corollary}
\newtheorem{fact}[theorem]{Fact}
\theoremstyle{definition}
\newtheorem{definition}[theorem]{Definition}
\newtheorem{conjecture}[theorem]{Conjecture}
\theoremstyle{remark}
\newtheorem{remark}[theorem]{Remark}

\newcommand{\tr}{\operatorname{tr}}
\newcommand{\Tr}{\operatorname{Tr}}
\newcommand{\sech}{\operatorname{sech}}
\newcommand{\frc}[1]{\{#1\}}

\title[The upper bound in the Collins--Shor conjecture]{The upper bound in the
Collins--Shor conjecture on the coefficients of the distance characteristic
polynomial of a tree\\[0.5ex] \normalfont\small (unreviewed draft prepared with AI assistance)}
\author[Conjectura project (draft)]{Conjectura project (authorship not yet determined; see the attribution section)}
\address{Conjectura project, directed by Gustavo Bernardo Brito, Universidade
Federal de Goi\'as (UFG), Brazil}
\date{Draft of 27 November 2026}

\emergencystretch=3em
\begin{document}

\begin{abstract}
For a tree $T$ of order $n\ge 3$ with distance matrix $D$, let
$d_0,\dots,d_{n-2}$ be the normalized coefficients
$d_k=2^{k-(n-2)}|\delta_k|$, where $\delta_k$ is the coefficient of $x^k$ in
$\det(D-xI)$. We give a written candidate proof, checked so far only by AI
reviews, that every index $p$ with
$d_p=\max_k d_k$ satisfies $p<n-n/\sqrt5+3/5$, and in particular
$p\le\lceil n-n/\sqrt5\rceil$. This is the upper half of the conjecture of
Collins and Shor in the form stated as Conjecture~1.2 by Aalipour et al.,
under the reading that every maximizing index is to be bounded. The proof
writes the polynomial as $4(\alpha-x)q(x)$ with $q$ real-rooted, bounds the
peak by the mode of $q$, the mode by the mean (Darroch's rule), and the mean
by $n-\tr(I+L)^{-1}$ through interlacing with the Laplacian spectrum; it then
shows that the path minimizes $\tr(I+L)^{-1}$ among trees of order $n$, and
evaluates this trace for the path. The minimality of the path is very
probably a consequence of a published theorem of Andriantiana,
Razanajatovo Misanantenaina and Wagner; the proof given here is independent
of theirs. Finally, from Collins's closed formula for the path we locate the
peak of the path exactly, which shows that the bound
$\lceil n-n/\sqrt5\rceil$ is attained for infinitely many $n$. No claim of
novelty is made; the searches made are listed.
\end{abstract}

\maketitle

\begin{center}
\fbox{\parbox{0.93\textwidth}{\small
\textbf{Unreviewed draft prepared with AI assistance. Not for distribution.}
The mathematics of this draft (the proof strategy, and the lemmas, theorems
and proofs of Sections~\ref{sec:spectral} to~\ref{sec:path}) was produced by
an AI system, Claude, model \nolinkurl{claude-opus-5-5}
(Anthropic), in the Conjectura project in November 2026; the text was
assembled from the project's notes by the same model. The proofs have been
checked only by AI reviews run in fresh contexts (models
\nolinkurl{claude-opus-5-5} and \nolinkurl{claude-sonnet-5-5}); \textbf{AI review is
not peer review}. \textbf{No mathematician has read the proof of either half
of the conjecture}, neither the upper half proved here nor the lower half,
for which the project holds a separate candidate proof. Novelty is not
claimed (Section~\ref{sec:literature}). Nothing in this draft has been
submitted, posted or sent to anyone; every outside use is decided by the
project director. The AI-use disclosure and the attribution record are at
the end.}}
\end{center}

\section{Introduction}\label{sec:intro}

Let $T$ be a tree of order $n$ and $D$ its distance matrix. By Graham and
Lov\'asz~\cite{GL78}, the coefficients of the characteristic polynomial of
$D$ have signs that allow them to be normalized to positive numbers
$d_0,\dots,d_{n-2}$ (Definition~\ref{def:dk}; \cite[Fact~1.1]{AAB}). The
conjecture of Graham and Lov\'asz, as Collins recalls it
\cite[p.~29]{Col89}, was that this sequence is unimodal with peak at
$k=\frac12n$; Collins showed that this is true for stars but not for paths,
computed the coefficients of the star and of the path, and stated as
Conjecture~9 of~\cite{Col89}, attributed there to Peter Shor:
``The coefficients of $\mathrm{CPD}(T)$ for any tree $T$ with $n$ vertices
are unimodal with peak between $\tfrac12 n$ and $n(1-1/\sqrt5)$.''
Aalipour, Abiad, Berikkyzy, Hogben, Kenter, Lin and Tait~\cite{AAB} proved
unimodality, indeed log-concavity, of $d_0,\dots,d_{n-2}$ for every tree
(\cite[Theorem~2.1]{AAB}), proved that the peak is at most
$\lceil 2n/3\rceil$ (\cite[Theorem~3.3]{AAB}), and restated the location of
the peak with a floor and a ceiling:

\begin{conjecture}[{\cite[Conjecture~1.2]{AAB}}; Collins--Shor]\label{conj:CS}
Verbatim, as transcribed in the project's statement file (the bracketed word
is part of that transcription): ``The [normalized] coefficients of the
distance characteristic polynomial for any tree $T$ with $n$ vertices are
unimodal with peak between $\lfloor n/2\rfloor$ and
$\lceil n-n/\sqrt5\rceil$.''
\end{conjecture}

The source does not say how a tie between two maximal coefficients is to be
read. Throughout we use the strong reading: \emph{every} index $p$ with
$d_p=\max_{0\le k\le n-2}d_k$ is to lie in
$[\lfloor n/2\rfloor,\lceil n-n/\sqrt5\rceil]$. Since unimodality is
\cite[Theorem~2.1]{AAB}, only the location of the peak is open. The
\emph{upper half} of Conjecture~\ref{conj:CS} is the statement that, for
every $n\ge3$ and every tree $T$ of order $n$, every maximizing index $p$
satisfies $p\le\lceil n-n/\sqrt5\rceil$. This draft gives a written proof of
it (Theorem~\ref{thm:T9}), with a slightly sharper bound.

\subsection*{Results}
For a graph $H$ let $L_H$ be its Laplacian and
$\Tr(H):=\tr(I+L_H)^{-1}$. Let $P_n$ be the path of order $n$.

\begin{itemize}
\item Theorem~\ref{thm:T9} (local name T9). For every tree $T$ of order
$n\ge3$, every maximizing index $p$ of $d_0,\dots,d_{n-2}$ satisfies
$p<n-n/\sqrt5+3/5$; in particular $p\le\lceil n-n/\sqrt5\rceil$.
\item Theorem~\ref{thm:T10} (T10). For every tree $T$ of order $n\ge1$,
$\Tr(T)\ge\Tr(P_n)$, with equality only for $T\cong P_n$.
\item Lemma~\ref{lem:S46} (S46). $\Tr(P_n)=n/\sqrt5+2/5+(2n/\sqrt5)\,r^{2n}/(1-r^{2n})$,
$r=\varphi^{-2}$, $\varphi=(1+\sqrt5)/2$.
\item Lemma~\ref{lem:S47} and Corollary~\ref{cor:S48} (S47, S48). The exact
peak of the path, and the attainment of $\lceil n-n/\sqrt5\rceil$ by the path
for infinitely many $n$.
\end{itemize}
The local names (S40 to S48, T9, T10) are those of the project notes, and
are kept so that each statement can be traced to its source file and its
reviews.

The lower half of Conjecture~\ref{conj:CS} (every maximizing index is at
least $\lfloor n/2\rfloor$) has a separate candidate proof in the project
(Theorem~T6(c) of the note \nolinkurl{notes/lower-bound.md}, registry record
MD-0019). It is not reproduced here and nothing in this draft depends on it,
except Corollary~\ref{cor:both}, which is stated conditionally. Neither half
has been read by a mathematician.

The proof of Theorem~\ref{thm:T9} uses only the classical facts collected in
Section~\ref{sec:tools} and the identification of the $d_k$ with the
coefficients of a normalized polynomial (Lemma~\ref{lem:ident}); it does not
use the log-concavity of $(d_k)$ \cite[Theorem~2.1]{AAB}.

\section{Definitions and the normalized polynomial}\label{sec:defs}

\begin{definition}\label{def:dk}
Let $T$ be a tree of order $n\ge2$ with distance matrix $D$, Laplacian $L$,
and let $\tau\in\mathbb R^{V(T)}$, $\tau_v=2-\deg v$. Let $\delta_k$ be the
coefficient of $x^k$ in $\det(D-xI)$. The \emph{normalized coefficients}
are $d_k=2^{k-(n-2)}|\delta_k|$, $0\le k\le n-2$. Put
\[
\ell_T(x)=-2^{-(n-2)}\det(2xI-D).
\]
A \emph{maximizing index} is an index $p\in\{0,\dots,n-2\}$ with
$d_p=\max_{0\le k\le n-2}d_k$. Put $M=\lceil n-n/\sqrt5\rceil$ and
$\varphi=(1+\sqrt5)/2$.
\end{definition}

\begin{lemma}[identification; {\cite[Observation~3.1]{AAB}}]\label{lem:ident}
For $0\le k\le n-2$ the coefficient of $x^k$ in $\ell_T$ is $d_k$.
\end{lemma}

\begin{proof}
$\det(2xI-D)=(-1)^n\det(D-2xI)=(-1)^n\sum_k\delta_k2^kx^k$, so the coefficient
of $x^k$ in $\ell_T$ is $(-1)^{n-1}2^{k-(n-2)}\delta_k$. By Graham and
Lov\'asz~\cite{GL78}, in the form \cite[Fact~1.1]{AAB}, $(-1)^{n-1}\delta_k>0$
for $0\le k\le n-2$, so this coefficient equals $2^{k-(n-2)}|\delta_k|=d_k$.
\end{proof}

For a graph $H$ write $G_H=(I+L_H)^{-1}$ and $\Tr(H)=\tr G_H$. A
\emph{rooted tree} $(X,v)$ is a tree $X$ with a distinguished vertex $v$; put
\[
g_X=(G_X)_{vv},\quad k_X=(G_X^2)_{vv},\quad \varepsilon_X=1/g_X-1,\quad
\kappa_X=k_X/g_X^2,\quad t_X=\tr G_X .
\]

\section{Classical tools}\label{sec:tools}

The following facts are used without proof. For each we record where the
project checked it.

\begin{enumerate}
\item[(GL)] \emph{Graham--Lov\'asz inverse.} For a tree of order $n\ge2$,
$D^{-1}=-\tfrac12L+\tau\tau^{\mathsf T}/(2(n-1))$ \cite[Lemma~1, p.~66]{GL78}.
Read in a scan of~\cite{GL78} by the project's literature audit (two minus
signs are illegible in the scan and were fixed from Fact~3 and eq.~(10) of
the same paper); verified in exact arithmetic by an AI reviewer for every
tree of order $3$ to $17$.
\item[(GP)] \emph{Graham--Pollak determinant.}
$\det D=(-1)^{n-1}(n-1)2^{n-2}$ \cite{GP71}. Not read at the source in the
project; it is quoted in \cite{AAB} (as $d_0=n-1$) and is implied by the
exact characteristic polynomials computed by the reviewers for all trees of
order at most $17$.
\item[(W)] \emph{Weyl interlacing for a rank-one positive semidefinite
update.} If $B=A+ww^{\mathsf T}$ with $A$ real symmetric, then, with
eigenvalues in decreasing order, $\lambda_i(A)\le\lambda_i(B)\le\lambda_{i-1}(A)$
for $2\le i\le n$, and $\lambda_1(A)\le\lambda_1(B)$. Textbook fact (e.g.\
\cite{HJ}); no specific statement in a source was checked in the project.
\item[(N)] \emph{Newton's inequalities.} If $\sum_{k=0}^Nq_kx^k$ has only real
zeros and positive coefficients, then
$q_k^2\ge q_{k-1}q_{k+1}(1+1/k)(1+1/(N-k))>q_{k-1}q_{k+1}$ for $1\le k\le N-1$.
Textbook fact (e.g.\ \cite{HLP}); not checked against a source in the
project.
\item[(D)] \emph{Darroch's rule} \cite[Theorem~4]{Dar64}, as stated by Pitman
\cite[p.~284]{Pit97}: if a polynomial $P$ with non-negative coefficients has
only real zeros, every index $m$ at which its coefficient is maximal
satisfies $|m-\mu|<1$, where $\mu=P'(1)/P(1)$. Darroch's paper was not read
in the project; the statement was read in Pitman~\cite{Pit97} by the
project's literature audit and, independently, by one AI reviewer, including
the case of two maximizing indices (\cite[eq.~(13)]{Pit97}). A second
reviewer, who could read neither source, confirmed the rule in a secondary
source and tested it on $20{,}000$ random real-rooted polynomials.
\item[(M)] \emph{M-matrix facts.} For a graph $H$, $I+L_H$ is a non-singular
M-matrix; if $H$ is connected, $(I+L_H)^{-1}$ is entrywise positive; and
$(I+L_H)\mathbf 1=\mathbf 1$, so every row of $G_H$ sums to $1$. Textbook
facts (e.g.\ \cite{BP}); not checked against a source in the project.
\item[(PS)] \emph{Spectrum of the path.} The Laplacian eigenvalues of $P_n$
are $2-2\cos(j\pi/n)$, $0\le j\le n-1$. Textbook fact (e.g.\ \cite{BH});
confirmed indirectly by the exact checks of Lemma~\ref{lem:S46}.
\item[(C)] \emph{Collins's formula for the path} \cite[Theorem~3 and its
proof, p.~31]{Col89}: for $n\ge3$,
$d_{n-k}(P_n)=\binom{n+k-1}{2k-1}\frac{(k-1)n}{k}$ for $2\le k\le n$. Read
first-hand in~\cite{Col89} by the project (the final step of Collins's proof
of the location of the peak is informal, and only the formula is used here);
checked in exact arithmetic against the distance matrix for $3\le n\le30$ by
the stored computation (Section~\ref{sec:comp}) and for $3\le n\le20$ by an
AI reviewer. Used only in Section~\ref{sec:path}.
\item[(K)] \emph{Kronecker's theorem}: for irrational $\xi$ the fractional
parts of $n\xi$, $n\ge1$, are dense in $[0,1)$ (indeed equidistributed, by
Weyl). Textbook fact. Used only in Corollary~\ref{cor:S48}.
\end{enumerate}
The Schur complement formula, the Sherman--Morrison formula and the
expansion of the Poisson kernel are used as elementary identities.

\section{The spectral form of \texorpdfstring{$\ell_T$}{l\_T} and the mode of \texorpdfstring{$q$}{q}}\label{sec:spectral}

\begin{lemma}[S40]\label{lem:S40}
Let $T$ be a tree of order $n\ge2$. The matrix $D/2$ has one positive
eigenvalue $\alpha$ and $n-1$ negative eigenvalues
$-\beta_1,\dots,-\beta_{n-1}$, and $\alpha=\sum_i\beta_i$. With
$q(x)=\prod_{i=1}^{n-1}(x+\beta_i)=\sum_{k=0}^{n-1}q_kx^k$,
\[
\ell_T(x)=4(\alpha-x)q(x),\qquad\text{so}\qquad d_k=4(\alpha q_k-q_{k-1})
\quad(0\le k\le n-2,\ q_{-1}=0).
\]
\end{lemma}

\begin{proof}
By (GL), $2D^{-1}=-L+ww^{\mathsf T}$ with $w=\tau/\sqrt{n-1}$. Since a tree is
connected, the eigenvalues of $-L$ in decreasing order are
$0\ge-\mu_{(1)}\ge\dots\ge-\mu_{(n-1)}$, where
$0<\mu_{(1)}\le\dots\le\mu_{(n-1)}$ are the non-zero Laplacian eigenvalues;
thus $\lambda_1(-L)=0$ and $\lambda_i(-L)=-\mu_{(i-1)}$ for $i\ge2$. Let
$\nu_1\ge\dots\ge\nu_n$ be the eigenvalues of $2D^{-1}$. By (W),
$\nu_i\le\lambda_{i-1}(-L)\le0$ for $i\ge2$; none is zero, since $D$ is
invertible by (GP); so $\nu_2,\dots,\nu_n<0$. The product $\prod_i\nu_i=\det(2D^{-1})$
has the sign of $\det D$, that is $(-1)^{n-1}$ by (GP); hence $\nu_1>0$. The
eigenvalues of $D/2$ are the $1/\nu_i$: $\alpha=1/\nu_1>0$ and, for $i\ge2$,
$-\beta_{i-1}=1/\nu_i<0$. Since $\tr D=0$, $\alpha=\sum_i\beta_i$. Finally,
with $\theta_1,\dots,\theta_n$ the eigenvalues of $D/2$,
$\det(2xI-D)=2^n\prod_i(x-\theta_i)=2^n(x-\alpha)q(x)$, so
$\ell_T=-2^{-(n-2)}2^n(x-\alpha)q(x)=4(\alpha-x)q(x)$. Comparing coefficients
and using Lemma~\ref{lem:ident} gives $d_k=4(\alpha q_k-q_{k-1})$ for
$0\le k\le n-2$.
\end{proof}

\begin{lemma}[S41: the peak of $d$ lies at or before a descent of $q$]\label{lem:S41}
Let $T$ be a tree of order $n\ge3$, with $\alpha$, $q$ as in
Lemma~\ref{lem:S40}, and let $0\le k\le n-3$. If $q_{k+1}<q_k$, then
$d_{j+1}<d_j$ for every $j$ with $k\le j\le n-3$; hence every maximizing
index of $d_0,\dots,d_{n-2}$ is at most $k$.
\end{lemma}

\begin{proof}
The polynomial $q$ has degree $n-1$, positive coefficients (its zeros
$-\beta_i$ are negative) and only real zeros. By (N),
$q_j^2>q_{j-1}q_{j+1}$ for $1\le j\le n-2$, so the ratios
$\varrho_j=q_{j+1}/q_j$, $0\le j\le n-2$, are strictly decreasing. Since $q$
is monic, $q_{n-2}=\sum_i\beta_i=\alpha$, so $\varrho_{n-2}=1/\alpha$, and
hence $\alpha\varrho_j>1$ for $0\le j\le n-3$. Fix $j$ with $k\le j\le n-3$;
then $\varrho_j\le\varrho_k<1$. By Lemma~\ref{lem:S40},
\[
d_{j+1}-d_j=4\bigl[\alpha(q_{j+1}-q_j)-(q_j-q_{j-1})\bigr].
\]
If $j=0$, this is $4[\alpha(q_1-q_0)-q_0]<0$. If $j\ge1$ and
$q_j\ge q_{j-1}$, the first term is negative and the second is non-positive.
If $j\ge1$ and $q_j<q_{j-1}$, then $\varrho_j<\varrho_{j-1}<1$, and writing
$q_j=\varrho_{j-1}q_{j-1}$, $q_{j+1}=\varrho_j\varrho_{j-1}q_{j-1}$,
\[
d_{j+1}-d_j=4q_{j-1}\bigl[(1-\varrho_{j-1})-\alpha\varrho_{j-1}(1-\varrho_j)\bigr];
\]
since $\alpha\varrho_{j-1}>1$ and $1-\varrho_j>1-\varrho_{j-1}>0$, the bracket
is negative. So $d_k>d_{k+1}>\dots>d_{n-2}$, and no index larger than $k$ is
maximizing.
\end{proof}

\begin{lemma}[S42: the mean of $q$ and the Laplacian spectrum]\label{lem:S42}
Let $T$ be a tree of order $n\ge2$, with $q$ as in Lemma~\ref{lem:S40}. Then
\[
\mu_q:=\frac{q'(1)}{q(1)}=\sum_{i=1}^{n-1}\frac1{1+\beta_i}
\le\sum_{j=1}^{n-1}\frac{\mu_{(j)}}{1+\mu_{(j)}}=n-\Tr(T).
\]
\end{lemma}

\begin{proof}
The first equality is the logarithmic derivative of $\prod_i(x+\beta_i)$ at
$x=1$. In the notation of the proof of Lemma~\ref{lem:S40},
$\beta_{i-1}=-1/\nu_i=1/|\nu_i|$ for $2\le i\le n$, so
$1/(1+\beta_{i-1})=|\nu_i|/(1+|\nu_i|)$. By (W), $\nu_i\ge\lambda_i(-L)=-\mu_{(i-1)}$
for $i\ge2$, so $|\nu_i|\le\mu_{(i-1)}$; and $y\mapsto y/(1+y)$ is increasing
on $(0,\infty)$. This gives the inequality. Finally
$\Tr(T)=1+\sum_j1/(1+\mu_{(j)})$ (the eigenvalue $0$ of $L$ contributes $1$),
so $\sum_j\mu_{(j)}/(1+\mu_{(j)})=\sum_j\bigl(1-1/(1+\mu_{(j)})\bigr)=(n-1)-(\Tr(T)-1)=n-\Tr(T)$.
\end{proof}

\begin{remark}
The interlacing between the distance eigenvalues and the reciprocals of the
non-zero Laplacian eigenvalues used in Lemmas~\ref{lem:S40}
and~\ref{lem:S42} is a theorem of Merris~\cite{Mer90}:
$0>-2/\mu_{(n-1)}\ge\partial_2\ge-2/\mu_{(n-2)}\ge\dots\ge\partial_n\ge-2/\mu_{(1)}$,
where $\partial_1\ge\dots\ge\partial_n$ are the eigenvalues of $D$; the
inertia of $D$ is $(1,n-1,0)$. Merris's paper was read in full by the
project's literature audit; the equation number cited for this chain in the
project notes was not confirmed by a reviewer. The proofs above are the
same argument, through (GL).
\end{remark}

\section{The trace of \texorpdfstring{$(I+L)^{-1}$}{(I+L)\^{}(-1)}}\label{sec:trace}

\begin{lemma}[S43: gluing two rooted trees]\label{lem:S43}
Let $(X,v)$ and $(Y,v)$ be rooted trees and let $T$ be the tree obtained by
identifying their roots. Then $1/g_T=1/g_X+1/g_Y-1$ (so
$\varepsilon_T=\varepsilon_X+\varepsilon_Y$), and
\[
t_T=t_X+t_Y-g_T\bigl[1+\varepsilon_Y\kappa_Xg_X+\varepsilon_X\kappa_Yg_Y\bigr].
\]
\end{lemma}

\begin{proof}
If $Y=\{v\}$ then $g_Y=1$, $\varepsilon_Y=0$, $t_Y=1$, $T=X$, and both claims
are immediate; similarly if $X=\{v\}$. Otherwise order the vertices of $T$ as
the vertices of $X$ (including $v$), then those of $Y\setminus v$. Let
$C=(I+L_Y)[Y\setminus v]$ be the principal submatrix of $I+L_Y$ on
$Y\setminus v$ (positive definite, as a principal submatrix of the positive
definite matrix $I+L_Y$), let $b\in\mathbb R^{Y\setminus v}$ be the indicator
vector of the neighbours of $v$ in $Y$, and let $\tilde b$ be the
$|X|\times|Y\setminus v|$ matrix whose row $v$ is $b^{\mathsf T}$ and whose
other rows are zero. Then
\[
I+L_T=\begin{pmatrix}I+L_X+\deg_Y(v)E_{vv}&-\tilde b\\-\tilde b^{\mathsf T}&C\end{pmatrix}.
\]
By the Schur complement formula, the $X$-block of $G_T$ is
$\bigl(I+L_X+sE_{vv}\bigr)^{-1}$ with $s=\deg_Y(v)-b^{\mathsf T}C^{-1}b$
(because $\tilde bC^{-1}\tilde b^{\mathsf T}=(b^{\mathsf T}C^{-1}b)E_{vv}$).
In $Y$ alone, ordered with $v$ first,
$I+L_Y=\begin{pmatrix}1+\deg_Y(v)&-b^{\mathsf T}\\-b&C\end{pmatrix}$, and the
same elimination gives $1/g_Y=1+\deg_Y(v)-b^{\mathsf T}C^{-1}b$; hence
$s=\varepsilon_Y$. By the Sherman--Morrison formula,
\[
(I+L_X+sE_{vv})^{-1}=G_X-\frac{s\,G_Xe_ve_v^{\mathsf T}G_X}{1+s\,g_X};
\]
its $(v,v)$ entry is $g_X/(1+sg_X)$, and its trace is
$t_X-sk_X/(1+sg_X)$ (as $\tr(G_Xe_ve_v^{\mathsf T}G_X)=e_v^{\mathsf T}G_X^2e_v=k_X$).
Hence $g_T=g_X/(1+\varepsilon_Yg_X)$, that is $1/g_T=1/g_X+\varepsilon_Y$,
which is the first claim. Exchanging the roles of $X$ and $Y$, the
$Y$-block of $G_T$ (on the vertex set of $Y$, including $v$) has trace
$t_Y-\varepsilon_Xk_Y/(1+\varepsilon_Xg_Y)$. The two blocks are principal
submatrices of $G_T$ whose index sets share only $v$, so
\[
t_T=t_X-\frac{\varepsilon_Yk_X}{1+\varepsilon_Yg_X}+t_Y-\frac{\varepsilon_Xk_Y}{1+\varepsilon_Xg_Y}-g_T .
\]
Finally
$\varepsilon_Yk_X/(1+\varepsilon_Yg_X)=\varepsilon_Y\kappa_Xg_X\cdot g_X/(1+\varepsilon_Yg_X)=g_T\,\varepsilon_Y\kappa_Xg_X$,
and symmetrically $\varepsilon_Xk_Y/(1+\varepsilon_Xg_Y)=g_T\,\varepsilon_X\kappa_Yg_Y$.
\end{proof}

\begin{lemma}[S43$'$: bounds at a root]\label{lem:S43p}
Let $(X,v)$ be a rooted tree with $|X|\ge2$. Then $0<g_X<1$, so
$a:=1/g_X=1+\varepsilon_X>1$, and $1\le\kappa_X\le1+\varepsilon_X/2=(1+a)/2$.
\end{lemma}

\begin{proof}
Write $G=G_X$. By (M), row $v$ of $G$ is positive and sums to $1$; as
$|X|\ge2$, $0<G_{vv}<1$. Hence
$\varepsilon_X=(1-G_{vv})/G_{vv}=\sum_{w\ne v}G_{vw}/G_{vv}>0$, and, $G$ being
symmetric, $\kappa_X=\sum_wG_{vw}^2/G_{vv}^2=1+\sum_{w\ne v}(G_{vw}/G_{vv})^2\ge1$.
For the upper bound it is enough to show $G_{wv}\le G_{vv}/2$ for every
$w\ne v$, since then
$\sum_{w\ne v}(G_{vw}/G_{vv})^2\le\frac12\sum_{w\ne v}G_{vw}/G_{vv}=\varepsilon_X/2$.
Let $w^*$ maximize $G_{wv}$ over $w\ne v$, and let $d=\deg w^*\ge1$. Row
$w^*$ of $(I+L)G=I$, in column $v$, reads $(1+d)G_{w^*v}=\sum_{x\sim w^*}G_{xv}$.
If $v$ is not a neighbour of $w^*$, the right side is at most $dG_{w^*v}$,
so $G_{w^*v}\le0$, contradicting (M). So $v\sim w^*$, and
$(1+d)G_{w^*v}\le G_{vv}+(d-1)G_{w^*v}$, that is $2G_{w^*v}\le G_{vv}$.
\end{proof}

For $N\ge2$ let $P=P_{N+1}$ with vertices $0,1,\dots,N$ in order, $G=G_P$,
and for a vertex $r$ put $E_r=1/G_{rr}-1$ and $c_r=(G^2)_{rr}/G_{rr}$.

\begin{lemma}[S44: the path seen from one vertex]\label{lem:S44}
Let $N\ge2$ and $1\le r\le N-1$. Then (i) $E_r>E_0$; (ii) $c_r>3/5$;
(iii) $c_0\ge c_r$, with equality only for $N=2$, $r=1$.
\end{lemma}

\begin{proof}
Let $G(t)=(tI+L_P)^{-1}$ for $t>0$, and write $t=4\sinh^2u$ with $u>0$. Then
$t=1$ corresponds to $u=\log\varphi$, since
$\sinh(\log\varphi)=(\varphi-\varphi^{-1})/2=1/2$. Let $e_j(t)$ be the
determinant of the $j\times j$ tridiagonal matrix with diagonal
$(t+1,t+2,\dots,t+2)$ and entries $-1$ next to the diagonal ($e_0=1$); it is
the principal minor of $tI+L_P$ on a segment of $j$ consecutive vertices
that contains an end of $P$ and does not contain $r$. Expanding along the
last row, $e_j=(t+2)e_{j-1}-e_{j-2}$ for $j\ge2$. Since $t+2=2\cosh2u$ and
$e_1=t+1=2\cosh2u-1=\cosh3u/\cosh u$, and since
$2\cosh2u\cosh((2j-1)u)=\cosh((2j+1)u)+\cosh((2j-3)u)$, induction gives
\[
e_j=\cosh((2j+1)u)/\cosh u\qquad(j\ge0).
\]
For $m\ge2$, expanding $\det(tI+L_{P_m})$ at an end vertex gives
$(t+1)e_{m-1}-e_{m-2}$; with $t+1=2\cosh2u-1$ and
$2\cosh2u\cosh((2m-1)u)=\cosh((2m+1)u)+\cosh((2m-3)u)$ this equals
\[
\frac{\cosh((2m+1)u)-\cosh((2m-1)u)}{\cosh u}=\frac{2\sinh u\,\sinh(2mu)}{\cosh u};
\]
for $m=1$ the right side is $4\sinh^2u=t=\det(tI+L_{P_1})$. Removing $r$
splits $P$ into segments of $r$ and $s=N-r$ vertices. With $a=2r+1$,
$b=2s+1$, $a+b=2N+2$, Cramer's rule gives, for every $0\le r\le N$,
\begin{equation}\label{eq:Grr}
G_{rr}(t)=\frac{e_re_s}{\det(tI+L_P)}=\frac{\cosh(au)\cosh(bu)}{2\sinh u\cosh u\,\sinh((a+b)u)}.
\end{equation}

(i) By \eqref{eq:Grr}, for $r=0$ one has $a=1$ and $b=a+b-1$ with the same
$a+b$; so
\[
\frac{G_{rr}}{G_{00}}=\frac{\cosh(au)\cosh(bu)}{\cosh(u)\cosh((a+b-1)u)}
=\frac{\cosh((a+b)u)+\cosh((a-b)u)}{\cosh((a+b)u)+\cosh((a+b-2)u)}<1,
\]
because $|a-b|<a+b-2$ when $a,b\ge3$, that is when $1\le r\le N-1$. So
$G_{rr}<G_{00}$ at $t=1$, and $E_r>E_0$.

(ii) Since $\frac{d}{dt}G(t)=-G(t)^2$, $c_r=-\frac{d}{dt}\log G_{rr}(t)$ at
$t=1$. From \eqref{eq:Grr},
\[
\frac{d}{du}\log G_{rr}=a\tanh(au)+b\tanh(bu)-\tanh u-\coth u-(a+b)\coth((a+b)u),
\]
and $dt/du=4\sinh2u$. At $u=\log\varphi$: $\cosh u=\sqrt5/2$,
$\sinh2u=(\varphi^2-\varphi^{-2})/2=\sqrt5/2$, $\tanh u=1/\sqrt5$,
$\coth u=\sqrt5$. Hence
\begin{equation}\label{eq:cr}
c_r=\frac{6/\sqrt5+(a+b)\coth((a+b)u)-a\tanh(au)-b\tanh(bu)}{2\sqrt5},
\end{equation}
valid for every $0\le r\le N$ (for $r=0$, $a=1$ and $b=2N+1$). Since
$\coth>1>\tanh$ on $(0,\infty)$, $(a+b)\coth((a+b)u)>a+b>a\tanh(au)+b\tanh(bu)$,
so the numerator exceeds $6/\sqrt5$ and $c_r>3/5$.

(iii) In \eqref{eq:cr} the $\coth$ terms are the same for $r$ and for $0$.
Using $\tanh z=1-2/(e^{2z}+1)$, and noting that the linear parts cancel
because $a+b=1+(2N+1)$, one gets at $u=\log\varphi$ (where
$e^{2xu}=\varphi^{2x}$)
\[
c_0-c_r=\frac{\psi(1)+\psi(2N+1)-\psi(a)-\psi(b)}{\sqrt5},\qquad
\psi(x)=\frac{x}{\varphi^{2x}+1}.
\]
Put $f_k=\psi(2k+1)$ and $\delta_k=f_k-f_{k+1}$. Then $\psi(a)=f_r$,
$\psi(b)=f_s$, and, as $N=s+r$,
\[
f_0+f_N-f_r-f_s=(f_0-f_r)-(f_s-f_{s+r})=\sum_{i=0}^{r-1}(\delta_i-\delta_{s+i}).
\]
It suffices to show that $\delta$ is non-increasing with $\delta_0=\delta_1$
and $\delta_k>\delta_{k+1}$ for $k\ge1$. Then every term is $\ge0$ (as
$s+i>i$). If $r\ge2$, the term $i=1$ is $\delta_1-\delta_{s+1}>0$; if $r=1$,
the only term is $\delta_0-\delta_s$, which vanishes exactly when $s=1$. So
equality holds only for $r=s=1$, that is $N=2$, $r=1$.

\emph{$\delta_0=\delta_1$.} For odd $x$,
$\varphi^{2x}+1=\varphi^x(\varphi^x+\varphi^{-x})=\sqrt5F_x\varphi^x$, where
$F_x$ is the Fibonacci number (Binet's formula). With $F_1=1$, $F_3=2$,
$F_5=5$,
\[
\psi(1)-2\psi(3)+\psi(5)=\frac{\varphi^{-1}-3\varphi^{-3}+\varphi^{-5}}{\sqrt5}
=\frac{\varphi^{-5}(\varphi^4-3\varphi^2+1)}{\sqrt5}=0,
\]
since $\varphi^4=3\varphi+2$ and $3\varphi^2=3\varphi+3$.

\emph{$\delta_k>\delta_{k+1}$ for $k\ge1$.} Since
$\delta_k-\delta_{k+1}=\psi(2k+1)-2\psi(2k+3)+\psi(2k+5)$, this follows from
strict convexity of $\psi$ on $[3,\infty)$. With $\gamma=2\log\varphi$ and
$\sigma(x)=1/(e^{\gamma x}+1)$ one has $\psi=x\sigma$,
$\sigma'=-\gamma\sigma(1-\sigma)$, and
\[
\psi''(x)=\gamma\sigma(1-\sigma)\bigl[\gamma x\tanh(\gamma x/2)-2\bigr],
\]
which is positive if and only if $z\tanh z>1$ for $z=x\log\varphi$. The
function $z\tanh z$ is increasing on $(0,\infty)$. At $x=3$,
$\tanh(3\log\varphi)=(\varphi^6-1)/(\varphi^6+1)=(8\varphi+4)/(8\varphi+6)=2/\sqrt5>0.894$
and $3\log\varphi>1.443$, so $z\tanh z>1.29$ for every $x\ge3$.
\end{proof}

\begin{lemma}[S45: merging two pendant paths lowers the trace]\label{lem:S45}
Let $(X,v)$ be a rooted tree with $|X|\ge2$ and let $j_1,j_2\ge1$. Let
$T(j_1,j_2)$ be the tree obtained from $X$ by attaching at $v$ two pendant
paths with $j_1$ and $j_2$ new vertices, and $T(j_1+j_2,0)$ the tree
obtained from $X$ by attaching at $v$ one pendant path with $j_1+j_2$ new
vertices. Then $\Tr(T(j_1,j_2))>\Tr(T(j_1+j_2,0))$.
\end{lemma}

\begin{proof}
Let $N=j_1+j_2\ge2$ and $P=P_{N+1}$ as before Lemma~\ref{lem:S44}. Then
$T(j_1,j_2)$ is $X$ glued at $v$ to $(P,j_1)$, and $T(N,0)$ is $X$ glued at
$v$ to $(P,0)$; note $1\le j_1\le N-1$. In Lemma~\ref{lem:S43} with
$(Y,v)=(P,r)$ one has $\varepsilon_Y=E_r$, $\kappa_Yg_Y=k_Y/g_Y=c_r$, and
$t_Y=\Tr(P)$, which does not depend on $r$. Put $a=1/g_X$ and
$\kappa=\kappa_X$; then $\varepsilon_X=a-1$ and $g_T=1/(a+E_r)$, so
Lemma~\ref{lem:S43} gives
\begin{gather*}
\Tr(T(N,0))=t_X+\Tr(P)-\Theta(E_0,c_0),\\
\Tr(T(j_1,j_2))=t_X+\Tr(P)-\Theta(E_{j_1},c_{j_1}),\\
\Theta(E,c)=\frac{1+E\kappa/a+(a-1)c}{a+E}.
\end{gather*}
Write $E=E_{j_1}$, $c=c_{j_1}$, $E'=E_0$, $c'=c_0$. Then
$\Tr(T(j_1,j_2))-\Tr(T(N,0))=\Theta(E',c')-\Theta(E,c)$. Multiplying
$\Theta(E',c')$ by $a(a+E')$ gives $a[1+(a-1)c']+\kappa E'$, and similarly
for $\Theta(E,c)$; expanding, the terms in $\kappa$ give $\kappa a(E'-E)$, and
\begin{gather*}
a(a+E)(a+E')\bigl[\Theta(E',c')-\Theta(E,c)\bigr]=a\,\mathcal N,\\
\mathcal N=(E-E')(1-\kappa)+(a-1)\bigl[(a+E)c'-(a+E')c\bigr].
\end{gather*}
By Lemma~\ref{lem:S44}(i), $E>E'$, and by Lemma~\ref{lem:S43p},
$1-\kappa\ge-(a-1)/2$; so
$\mathcal N\ge(a-1)\bigl[(a+E)c'-(a+E')c-(E-E')/2\bigr]$. By
Lemma~\ref{lem:S44}(iii), $c'\ge c$, so $(a+E)c'\ge(a+E)c$ and
\[
\mathcal N\ge(a-1)(E-E')(c-1/2),
\]
which is positive by Lemma~\ref{lem:S43p} ($a>1$), Lemma~\ref{lem:S44}(i)
and Lemma~\ref{lem:S44}(ii) ($c>3/5$). As $a(a+E)(a+E')>0$, the claim
follows.
\end{proof}

\begin{theorem}[T10: the path minimizes $\tr(I+L)^{-1}$]\label{thm:T10}
For every tree $T$ of order $n\ge1$, $\tr(I+L_T)^{-1}\ge\tr(I+L_{P_n})^{-1}$,
with equality only for $T\cong P_n$. Equivalently, among trees of order $n$,
$\sum_j\mu_j/(1+\mu_j)$ over the Laplacian eigenvalues $\mu_j$ is maximized
exactly by the path.
\end{theorem}

\begin{proof}
By induction on the number of leaves. A tree with at most two leaves is a
path. Otherwise $T$ has a vertex of degree at least $3$. Root $T$ at a leaf
$\rho$ and let $v$ be a vertex of degree $\ge3$ at maximum distance from
$\rho$. Every branch at $v$ (component of $T-v$) not containing $\rho$
contains no vertex of degree $\ge3$, since each of its vertices is farther
from $\rho$ than $v$; so it is a path, and its vertex adjacent to $v$ has
degree at most $2$ in $T$, hence is an end of that path. There are
$\deg v-1\ge2$ such branches. Take two of them, with $j_1$ and $j_2$
vertices, and let $X$ be $T$ without them; then $v\in X$ and $\rho\in X$
with $\rho\ne v$, so $|X|\ge2$, and $T=T(j_1,j_2)$ in the notation of
Lemma~\ref{lem:S45}. By Lemma~\ref{lem:S45}, $T'=T(j_1+j_2,0)$ has
$\Tr(T')<\Tr(T)$. The tree $T'$ has order $n$ and one leaf fewer than $T$:
the two ends of the two pendant paths are replaced by one, and $v$ keeps
degree $\deg v-1\ge2$. By induction $\Tr(T')\ge\Tr(P_n)$, so
$\Tr(T)>\Tr(P_n)$. The equivalent form follows from
$\sum_j\mu_j/(1+\mu_j)=n-\Tr(T)$ (proof of Lemma~\ref{lem:S42}).
\end{proof}

\begin{remark}[a second route; not reviewed]\label{rem:AMW}
The project's literature audit derived Theorem~\ref{thm:T10} from the work
of Andriantiana, Razanajatovo Misanantenaina and Wagner~\cite{AMW} as
follows. For a forest the matching polynomial is the characteristic
polynomial. The subdivision $S(T)$ of a tree $T$ is a tree whose adjacency
eigenvalues are $\pm\sqrt{\mu_j}$ (over the Laplacian eigenvalues $\mu_j$ of
$T$) together with zeros. By \cite[Lemma~5.1]{AMW}, the average size of a
matching satisfies $\mathrm{avm}(G)=\frac12\sum_\lambda\lambda^2/(\lambda^2+1)$
over the zeros $\lambda$ of the matching polynomial of $G$; hence
$\mathrm{avm}(S(T))=\sum_j\mu_j/(1+\mu_j)$. By \cite[Theorem~4.4]{AMW}, the
path maximizes $\mathrm{avm}$ among trees of a given order, strictly (its
proof uses computer-checked cases). Since $S(P_n)=P_{2n-1}$,
Theorem~\ref{thm:T10} would follow. This derivation is the auditor's; its two
classical steps (matching polynomial of a forest; spectrum of the
subdivision) were not checked against a source, and no reviewer checked the
derivation. If it is correct, Theorem~\ref{thm:T10} is a known consequence
of a published theorem, and the proof given above is an independent,
self-contained one. We regard Theorem~\ref{thm:T10} as very probably known.
\end{remark}

\begin{lemma}[S46: the value for the path]\label{lem:S46}
For $n\ge1$, with $r=\varphi^{-2}$,
\[
\tr(I+L_{P_n})^{-1}=\frac{n}{\sqrt5}+\frac25+\frac{2n}{\sqrt5}\,\frac{r^{2n}}{1-r^{2n}}>\frac{n}{\sqrt5}+\frac25 .
\]
\end{lemma}

\begin{proof}
By (PS), the eigenvalues of $I+L_{P_n}$ are $3-2\cos\theta_j$,
$\theta_j=j\pi/n$, $0\le j\le n-1$. Since $r^2-3r+1=0$ (equivalently
$\varphi^4-3\varphi^2+1=0$), $r+1/r=3$ and
\[
\frac1{3-2\cos\theta}=\frac{r}{1-2r\cos\theta+r^2}
=\frac1{\sqrt5}\Bigl(1+2\sum_{m\ge1}r^m\cos m\theta\Bigr),
\]
by the Poisson kernel, with $r/(1-r^2)=1/(\varphi^2-\varphi^{-2})=1/\sqrt5$;
the series converges absolutely, as $0<r<1$. Now
$\sum_{j=0}^{n-1}\cos(m\theta_j)=\operatorname{Re}\sum_{j=0}^{n-1}\omega^j$ with
$\omega=e^{i\pi m/n}$, and $\omega^n=e^{i\pi m}$. If $2n\mid m$ the sum is $n$.
If $m$ is even and $2n\nmid m$, then $\omega\ne1=\omega^n$ and the sum is $0$.
If $m$ is odd, then $m/n$ is not an even integer, so $\omega\ne1$, and
$\omega^n=-1$; the sum is
$\operatorname{Re}[(1-\omega^n)/(1-\omega)]=\operatorname{Re}[2/(1-\omega)]=1$,
because $\operatorname{Re}[1/(1-e^{i\psi})]=1/2$ whenever $e^{i\psi}\ne1$.
This covers $n=1$ (one eigenvalue, equal to $1$). Summing over $j$ and
interchanging the sums (absolute convergence),
\[
\tr(I+L_{P_n})^{-1}=\frac1{\sqrt5}\Bigl[n+2\sum_{m\text{ odd}}r^m+2n\sum_{l\ge1}r^{2nl}\Bigr],
\]
and $2\sum_{m\text{ odd}}r^m=2r/(1-r^2)=2/\sqrt5$, which gives the formula.
The last term is positive.
\end{proof}

\begin{remark}
The same value follows from \cite[Lemma~4.2, eq.~(8)]{AMW} applied to the
path with $2n-1$ vertices, as the literature audit noted (not checked by a
reviewer).
\end{remark}

\section{Proof of the upper bound}\label{sec:proof}

\begin{theorem}[T9: the upper half of Conjecture~\ref{conj:CS}]\label{thm:T9}
For every tree $T$ of order $n\ge3$, every maximizing index $p$ of the
normalized coefficients $d_0,\dots,d_{n-2}$ satisfies $p<n-n/\sqrt5+3/5$. In
particular $p\le\lceil n-n/\sqrt5\rceil$; more precisely
$p\le K:=\lceil n-n/\sqrt5-2/5\rceil$.
\end{theorem}

\begin{proof}
Let $x=n-n/\sqrt5$, which is irrational since $\sqrt5$ is, and
$K=\lceil x-2/5\rceil$; note $K\le\lceil x\rceil=M$, and $K\ge2$ because
$x-2/5\ge3(1-1/\sqrt5)-2/5>1$ for $n\ge3$. By Lemma~\ref{lem:S42},
Theorem~\ref{thm:T10} and Lemma~\ref{lem:S46},
\[
\mu_q\le n-\Tr(T)\le n-\Tr(P_n)<x-2/5 .
\]
By (D), applied to $q$ (positive coefficients, only real zeros), every
maximizing index $m_q$ of $q_0,\dots,q_{n-1}$ satisfies
$m_q<\mu_q+1<(x-2/5)+1$. As $m_q$ is an integer and $x-2/5$ is not,
$m_q\le\lceil x-2/5\rceil=K$.

If $K\ge n-2$, every maximizing index $p$ of $d_0,\dots,d_{n-2}$ satisfies
$p\le n-2\le K$. Otherwise $0\le K\le n-3$, and $q_{K+1}<q_K$: if not, then,
the ratios $\varrho_j=q_{j+1}/q_j$ being strictly decreasing (proof of
Lemma~\ref{lem:S41}), $\varrho_j>\varrho_K\ge1$ for $j<K$, so
$q_0<q_1<\dots<q_K\le q_{K+1}$, and the maximum of $q_0,\dots,q_{n-1}$ would
be attained at some index $\ge K+1$, contradicting $m_q\le K$ for every
maximizing index $m_q$. Lemma~\ref{lem:S41} with $k=K$ then gives $p\le K$.
In both cases $p\le K\le M$ and $p\le K<(x-2/5)+1=x+3/5$.
\end{proof}

\begin{corollary}[conditional on the lower half]\label{cor:both}
Suppose that, for every tree $T$ of order $n\ge3$, every maximizing index of
$d_0(T),\dots,d_{n-2}(T)$ is at least $\lfloor n/2\rfloor$ (the lower half
of Conjecture~\ref{conj:CS}; the project's candidate proof of it is
Theorem~T6(c) of \nolinkurl{notes/lower-bound.md}, record MD-0019, not
reproduced or assessed here). Then Conjecture~\ref{conj:CS} holds for every
tree, in the strong reading: every maximizing index lies in
$[\lfloor n/2\rfloor,\lceil n-n/\sqrt5\rceil]$.
\end{corollary}

\begin{proof}
Combine the hypothesis with Theorem~\ref{thm:T9}; unimodality is
\cite[Theorem~2.1]{AAB}.
\end{proof}

Both halves are candidate proofs; neither is established, and neither has
been read by a mathematician.

\begin{corollary}\label{cor:sharper}
(a) For every tree of order $n\ge3$, every maximizing index is at most
$K=\lceil n-n/\sqrt5-2/5\rceil$, which equals $\lfloor n-n/\sqrt5\rfloor$
when the fractional part of $n-n/\sqrt5$ is below $2/5$. In particular, the
peak of the path can equal $\lceil n-n/\sqrt5\rceil$ only when that
fractional part exceeds $2/5$ (in fact only when it exceeds $4/5$, by
Corollary~\ref{cor:S48}).

(b) For an individual tree $T$ of order $n\ge3$, every maximizing index $p$
satisfies $p<1+\sum_j\mu_j/(1+\mu_j)=n+1-\tr(I+L_T)^{-1}$.
\end{corollary}

\begin{proof}
(a) The first statement is Theorem~\ref{thm:T9}. If the fractional part
$f$ of $x=n-n/\sqrt5$ is below $2/5$, then
$x-2/5\in(\lfloor x\rfloor-1,\lfloor x\rfloor)$, so $K=\lfloor x\rfloor$, which
is less than $\lceil x\rceil$ as $x$ is irrational. The last sentence follows
since the path is a tree, with Corollary~\ref{cor:S48} for the refinement.

(b) Let $k^*$ be the largest maximizing index of $q_0,\dots,q_{n-1}$. By (D)
and Lemma~\ref{lem:S42}, $k^*<\mu_q+1\le n+1-\Tr(T)$. If $k^*\le n-3$, then
$q_{k^*+1}<q_{k^*}$ (as $k^*+1$ is not maximizing), and Lemma~\ref{lem:S41}
with $k=k^*$ gives $p\le k^*$. Otherwise $p\le n-2\le k^*$. In both cases
$p\le k^*<n+1-\Tr(T)$.
\end{proof}

\begin{remark}
For the bound $p\le\lceil n-n/\sqrt5\rceil$ of the conjecture alone,
$\Tr(T)>n/\sqrt5$ is enough: it gives $\mu_q<x$, hence $m_q<x+1$, so
$m_q\le\lceil x\rceil$, and then Lemma~\ref{lem:S41} as in the proof of
Theorem~\ref{thm:T9} with $K$ replaced by $\lceil x\rceil$. The extra $2/5$
of Lemma~\ref{lem:S46} is used only for the sharper bound $K$.
Theorem~\ref{thm:T10} is stronger than either requirement. The statement
that the path has the largest peak index among all trees of order $n$ is
not proved here for general $n$, and Theorem~\ref{thm:T9} does not need it
(but see the end of Section~\ref{sec:path}).
\end{remark}

\section{The peak of the path}\label{sec:path}

In this section $n\ge3$, $d_j=d_j(P_n)$, $x=n-n/\sqrt5$ and $y=n/\sqrt5$
(both irrational). All statements of this section rest on Collins's formula
(C). For every integer $k$ put
\[
h_n(k)=5k^3+2k^2-4k-2-kn^2,
\]
and let $k_0(n)$ be the least integer $k\ge2$ with $h_n(k)>0$.

\begin{lemma}[S47: the peak of the path]\label{lem:S47}
Let $n\ge3$. Then $k_0(n)$ exists and $k_0(n)\le n-1$, and
$d_0,\dots,d_{n-2}$ is strictly increasing up to the index $n-k_0(n)$ and
strictly decreasing after it; so the path has the unique maximizing index
$p(P_n)=n-k_0(n)$. Moreover $p(P_n)\in\{\lfloor x\rfloor,\lceil x\rceil\}$, and
$p(P_n)=\lceil x\rceil$ if and only if $h_n(\lfloor y\rfloor)>0$.
\end{lemma}

\begin{proof}
For $2\le k\le n-1$, by (C) (applied to $k$ and to $k+1\le n$) and
$\binom{n+k}{2k+1}/\binom{n+k-1}{2k-1}=\frac{(n+k)(n-k)}{2k(2k+1)}$,
\[
\frac{d_{n-k}}{d_{n-k-1}}=\frac{k-1}{k}\cdot\frac{k+1}{k}\cdot\frac{2k(2k+1)}{(n+k)(n-k)}
=\frac{2(k^2-1)(2k+1)}{k(n^2-k^2)} .
\]
All factors are positive, and $2(k^2-1)(2k+1)-k(n^2-k^2)=h_n(k)$; so
$d_{n-k}>d_{n-k-1}$ if and only if $h_n(k)>0$, and likewise with $<$ and $=$.
Now $h_n(k)/k=5k^2+2k-4-2/k-n^2$ is strictly increasing in real $k>0$, and
$h_n(k)\ne0$ for every integer $k\ge2$: a zero would make $k$ divide $2$
(as $h_n(k)\equiv-2\pmod k$), so $k=2$, and $h_n(2)=38-2n^2\ne0$. Hence,
among integers $k\ge2$, $h_n(k)>0$ exactly for $k\ge k_0(n)$, once $k_0(n)$
exists. And $k_0(n)\le n-1$:
$h_n(n-1)=(n-1)\bigl(5(n-1)^2-n^2\bigr)+2(n-1)^2-4(n-1)-2$ is positive for
$n\ge4$ (both parts are, since $\sqrt5(n-1)>n$, and $2k^2-4k-2>0$ for
$k=n-1\ge3$), and equals $20$ for $n=3$. Reading $k=n-1,n-2,\dots,2$, that is
the index pairs $(n-k-1,n-k)$ from $(0,1)$ upwards, $d$ increases strictly
while $k\ge k_0$ and decreases strictly once $k<k_0$; the unique peak is at
$n-k_0$.

\emph{Location for $n\ge6$.} Then $y\ge6/\sqrt5>1+\sqrt2$ (since
$\sqrt5(1+\sqrt2)<5.4$). For $k=\lceil y\rceil>y$, monotonicity gives
$h_n(k)/k>5y^2+2y-4-2/y-n^2=2y-4-2/y>0$ (as $y^2-2y-1>0$), so
$k_0\le\lceil y\rceil$. If $k=\lfloor y\rfloor-1\ge2$, then $k<y-1$ and
$h_n(k)/k<5(y-1)^2+2(y-1)-4-2/(y-1)-n^2=-8y-1-2/(y-1)<0$, so $h_n<0$ at all
integers $2\le k\le\lfloor y\rfloor-1$ and $k_0\ge\lfloor y\rfloor$; this is
trivially true also when $\lfloor y\rfloor-1<2$. Hence
$k_0\in\{\lfloor y\rfloor,\lceil y\rceil\}$ and
$p=n-k_0\in\{\lfloor x\rfloor,\lceil x\rceil\}$, because
$n-\lceil y\rceil=\lfloor x\rfloor$ and $n-\lfloor y\rfloor=\lceil x\rceil$.
Since $\lfloor y\rfloor\ge2$, $p=\lceil x\rceil$ if and only if
$k_0=\lfloor y\rfloor$, if and only if $h_n(\lfloor y\rfloor)>0$.

\emph{The cases $n=3,4,5$, directly.} Here $x\approx1.66,\,2.21,\,2.76$ and
$\lfloor y\rfloor=1,1,2$. One has $h_3(2)=20$, $h_4(2)=6$, $h_5(2)=-12$,
$h_5(3)=64$, so $k_0=2,2,3$ and $p=1,2,2=\lfloor x\rfloor$. For the criterion
(the polynomial $h_n$ being defined at every integer):
$h_3(1)=-8$, $h_4(1)=-15$, $h_5(2)=-12$ are negative and $p=\lfloor x\rfloor\ne\lceil x\rceil$,
so ``$p=\lceil x\rceil$ if and only if $h_n(\lfloor y\rfloor)>0$'' holds, both
sides being false.
\end{proof}

\begin{corollary}[S48: when the path reaches the ceiling]\label{cor:S48}
Let $n\ge3$ and let $\theta$ be the fractional part of $y=n/\sqrt5$.
(a) If $\theta\ge1/5$, then $p(P_n)=\lfloor x\rfloor$. (b) If $\theta<1/10$ and
$n\ge30$, then $p(P_n)=\lceil x\rceil$. (c) In particular,
$p(P_n)=\lceil n-n/\sqrt5\rceil$ implies that the fractional part of
$n-n/\sqrt5$ exceeds $4/5$, and $p(P_n)=\lceil n-n/\sqrt5\rceil$ holds for
infinitely many $n$. Thus the upper bound $\lceil n-n/\sqrt5\rceil$ of
Conjecture~\ref{conj:CS} is attained, by the path, for infinitely many $n$,
and for those $n$ it cannot be replaced by $\lfloor n-n/\sqrt5\rfloor$.
\end{corollary}

\begin{proof}
For $n\le5$, Lemma~\ref{lem:S47} gives $p=\lfloor x\rfloor$ directly, so (a)
holds and (b), (c) need only $n\ge6$. Let $n\ge6$, so that
$k=\lfloor y\rfloor=y-\theta\ge2$. Then
\[
h_n(k)/k=2y(1-5\theta)+5\theta^2-2\theta-4-2/k .
\]
(a) If $\theta\ge1/5$, the first term is $\le0$, $5\theta^2-2\theta-4<0$ on
$[0,1)$, and $-2/k<0$; so $h_n(\lfloor y\rfloor)<0$ and Lemma~\ref{lem:S47}
gives $p=\lfloor x\rfloor$. (b) If $\theta<1/10$ and $n\ge30$, the first term
exceeds $y\ge30/\sqrt5>13$, while $|5\theta^2-2\theta-4-2/k|<6$; so
$h_n(\lfloor y\rfloor)>0$ and $p=\lceil x\rceil$. (c) The fractional part of
$x=n-y$ is $1-\theta$; if $p=\lceil x\rceil$ then $\theta<1/5$ by (a), so
$1-\theta>4/5$. Since $1/\sqrt5$ is irrational, (K) gives infinitely many
$n$ with $\theta<1/10$, hence infinitely many such $n\ge30$, and for them
$p(P_n)=\lceil x\rceil$ by (b). As $p(P_n)$ is the unique maximizing index
(Lemma~\ref{lem:S47}) and $\lceil x\rceil>\lfloor x\rfloor$, the ceiling cannot
be replaced by the floor for those $n$.
\end{proof}

\emph{Consistency with Theorem~\ref{thm:T9}.} Theorem~\ref{thm:T9} gives
$p(T)\le\lceil x-2/5\rceil$ for every tree. When the path is at $\lceil x\rceil$,
the fractional part of $x$ exceeds $4/5>2/5$, so
$\lceil x-2/5\rceil=\lceil x\rceil$: no contradiction. When the fractional
part of $x$ is below $2/5$ (which implies $\theta=1-\frc{x}>3/5\ge1/5$),
Theorem~\ref{thm:T9} gives $p(T)\le\lfloor x\rfloor$ for every tree $T$ of
order $n$, and Corollary~\ref{cor:S48}(a) gives $p(P_n)=\lfloor x\rfloor$; so
for those $n$ the path has the largest peak index among trees of order $n$.

\section{Computational evidence}\label{sec:comp}

The computations below support, but are not part of, the proofs. They were
not rerun for this draft.

\emph{Stored computation.} \nolinkurl{experiments/CJ-0001-collins-shor-peak/upper_bound.py},
committed at \texttt{3b95386} with the regression test
\nolinkurl{tests/test_cj0001_upper_bound.py}; result file
\nolinkurl{experiments/CJ-0001-collins-shor-peak/results/upper-bound.json}. Command
(from the repository root):
\begin{verbatim}
uv run python \
  experiments/CJ-0001-collins-shor-peak/upper_bound.py \
  --out \
  experiments/CJ-0001-collins-shor-peak/results/upper-bound.json
\end{verbatim}
The stored run (block \texttt{run} of the result file) was made at commit
\texttt{917d657} in a clean worktree, with Python 3.12.4, python-flint 0.9.0
and networkx 3.7, in 298.7 s, with 0 failures in every part. Arithmetic is
exact (integers, rationals) except where stated. It checks: Theorem~\ref{thm:T10}
on all trees of order $2$ to $20$; Lemma~\ref{lem:S46} (the inequality
exactly for $n\le200$; the closed form to 50 significant digits for
$n\le60$, high-precision numeric, not exact); the bound $p<n-n/\sqrt5+3/5$
of Theorem~\ref{thm:T9} on all trees of order $3$ to $18$ and, through
Collins's formula, on the path for $n\le2000$ (the formula itself is checked
against the distance matrix only for $n\le30$; for $31\le n\le2000$ these
path checks rest on (C)); Lemma~\ref{lem:S45} on all rooted trees with
$2\le|X|\le8$ and $j_1+j_2\le8$; Lemma~\ref{lem:S44}(iii) for $N\le60$; the
trace formula of Lemma~\ref{lem:S43} on all pairs of rooted trees of order
at most $6$; and, in floating point (exploratory only), Lemma~\ref{lem:S42}
on all trees of order $3$ to $14$. The result file also contains a part V
(trees of order at most $9$ computed by three methods, according to the
script's help text), not described in the notes. For the path, the run
finds 393 values of $n\le2000$ with $p(P_n)=\lceil x\rceil$; the smallest
fractional part of $x$ among them is $0.80128$.

\emph{Reviewers' own computations} (scratch, not stored). The two AI reviews
of the upper bound checked every intermediate inequality and
Theorem~\ref{thm:T9} in exact or rigorous arithmetic on all trees of order
$3$ to $17$ (respectively $3$ to $16$), on paths up to order $160$
(respectively $3000$, via Collins's formula), and on $198$ further trees of
orders $20$ to $70$ (respectively $162$ trees of orders $17$ to $60$; random
trees and special families); the review of Section~\ref{sec:path} checked
Collins's formula against the distance matrix for $n\le20$ and
Lemma~\ref{lem:S47} and Corollary~\ref{cor:S48} exactly for $3\le n\le700$.
No counterexample to any statement was found. Finite checks establish the
statements only in the range checked.

\section{Formal verification}\label{sec:formal}

Only the sequence part of Lemma~\ref{lem:S41} is formalized: Lean~4
theorems \nolinkurl{Conjectura.Research.CollinsShorPeak.upperD_succ_lt} and
\nolinkurl{Conjectura.Research.CollinsShorPeak.upperD_maxIndex_le} in
\nolinkurl{formalization/Conjectura/Research/CollinsShorPeak/UpperBoundPeak.lean},
as statements about real sequences in which positivity of $q$, the strict
decrease of the ratios $q_{j+1}/q_j$ and $q_{n-1}/q_{n-2}=1/\alpha$ are
hypotheses. Nothing about trees, distance matrices, Lemma~\ref{lem:S40},
Darroch's rule, the trace of $(I+L)^{-1}$, Theorem~\ref{thm:T9} or
Section~\ref{sec:path} is formalized. The project records a passing receipt
(\nolinkurl{verification/receipts/20261127T192616Z-e16cb626ca-dirty-passed.json};
Lean 4.36.0, Mathlib commit \texttt{fc4f59e5}) made at commit
\texttt{e16cb62} on a working tree that was not clean; the comparison of the
Lean statements with the note was made by the agent that wrote them and has
not been independently reviewed. The check was not rerun for this draft.

\section{Literature and novelty}\label{sec:literature}

\emph{No claim of novelty is made.} What was searched, as recorded in the
project's dossier for this problem (audit for session 13, 2026-11-27, by an
AI literature-auditor, model \nolinkurl{claude-opus-5-5}): 15 web queries
(English web and arXiv only); the five papers citing \cite{AAB} found
through the Semantic Scholar API (arXiv:2407.03309, 2206.07561, 2104.14516,
2103.00647, 1812.05734); \cite{AMW} (arXiv version, pp.~1--18; the journal
version was not read); \cite{Pit97} (p.~284); \cite{Mer90} (in full); and
several abstracts. Earlier audits of the project (sessions 2 to 11) read
\cite{AAB}, parts of \cite{GL78}, \cite{Col89} in full and parts of Collins's
thesis. Not searched: MathSciNet, zbMATH, Google Scholar, MathOverflow, and
the papers citing \cite{AMW}. The two AI reviewers made a few further web
searches of their own. An unsuccessful search does not establish
originality.

Findings within this coverage only.
(1) No prior proof of the upper half of Conjecture~\ref{conj:CS} was found;
the bound $\lceil 2n/3\rceil$ of \cite[Theorem~3.3]{AAB} is proved by ratio
inequalities, and no source linking the peak with Darroch's rule was found.
The ingredients (the spectral form, Darroch's rule, Merris's interlacing) are
known tools.
(2) Theorem~\ref{thm:T10} is very probably a known consequence of
\cite[Theorem~4.4]{AMW} (Remark~\ref{rem:AMW}); our proof by merging pendant
paths is independent of theirs. No direct statement in Laplacian,
rooted-forest, resolvent-energy or Kirchhoff-index language was found.
(3) Lemma~\ref{lem:S46} also follows from \cite[Lemma~4.2]{AMW}.
(4) The location of the peak of the path up to one is the content of
\cite[Theorem~3]{Col89}, whose last step is informal; the exact criterion of
Lemma~\ref{lem:S47} and Corollary~\ref{cor:S48} was not the subject of a
separate literature search, and its novelty is unknown.

\section*{Status of the contributions}

The project keeps four independent status fields for each registered
contribution (registry \nolinkurl{research/discoveries.json}, values at the time
of writing, uncommitted working tree of 2026-11-27):

\smallskip
\noindent\begin{tabular}{@{}p{0.22\textwidth}p{0.74\textwidth}@{}}
MD-0022 (Thm.~\ref{thm:T9}) & mathematical: candidate-proof; formal: partial (Lemma~\ref{lem:S41}, sequence part, only); novelty: preliminary-search-completed (no prior proof found within the coverage above); external review: ai-reviewed. Standing: original candidate, not established.\\[0.5ex]
MD-0023 (Thm.~\ref{thm:T10}) & mathematical: candidate-proof; formal: not-started; novelty: preliminary-search-completed; external review: ai-reviewed. Standing: very probably a known result (Remark~\ref{rem:AMW}); the proof here is independent; not established.\\[0.5ex]
MD-0024 (Lem.~\ref{lem:S47}, Cor.~\ref{cor:S48}) & mathematical: candidate-proof; formal: not-started; novelty: unknown; external review: ai-reviewed. Standing: refinement of \cite[Theorem~3]{Col89}; not established.
\end{tabular}
\smallskip

\noindent In short: written proofs, AI-reviewed, not human-reviewed, not
peer-reviewed, novelty not established, essentially not formalized.

\section*{Use of AI systems}

\emph{Systems.} Claude (Anthropic), models \nolinkurl{claude-opus-5-5} and
\nolinkurl{claude-sonnet-5-5}, run through Claude Code in the Conjectura
repository in November 2026. Earlier project files used here
(the identification Lemma~\ref{lem:ident} and the reading of \cite{Col89})
were produced by \nolinkurl{claude-fable-5-2}.

\emph{What the systems did.} \nolinkurl{claude-opus-5-5}, as principal
investigator: the proof strategy (spectral form, Darroch's rule, reduction
to $\tr(I+L)^{-1}$), the statements and proofs of Lemmas S40 to S47,
Theorems T9 and T10 and Corollary S48, and the stored computation.
\nolinkurl{claude-opus-5-5}, as literature auditor: located Darroch's rule
through \cite{Pit97}, \cite{Mer90} and \cite{AMW}, and derived the second
route of Remark~\ref{rem:AMW}. \nolinkurl{claude-opus-5-5}, as formalization
agent: the Lean statements and proofs of Section~\ref{sec:formal}.
\nolinkurl{claude-opus-5-5}, as publication agent: the assembly and wording of
this draft from the notes. Reviews: (1) \nolinkurl{claude-opus-5-5}, fresh
context (same model as the author, so independent by context only), of the
upper-bound note at commit \texttt{53421a6}; (2) \nolinkurl{claude-sonnet-5-5},
fresh context (different model, same vendor), of the same revision; (3)
\nolinkurl{claude-sonnet-5-5}, recheck of the changes of revision~1 at commit
\texttt{2037b18}; (4) \nolinkurl{claude-sonnet-5-5}, of the path note at
\texttt{2037b18}. Each returned ``no gap found within stated scope'', and
each listed defects of wording, provenance or citation, which are folded
into this draft. No review covered this \LaTeX{} text itself.

\emph{What was checked by a human.} Nothing. No mathematician has read the
proofs of this draft or the project's candidate proof of the lower half. The
project director set the research direction only.

\emph{What was verified computationally or formally.} See
Sections~\ref{sec:comp} and~\ref{sec:formal}: finite exact checks in the
ranges stated; Lean verification of one sequence lemma under its stated
hypotheses.

\emph{What was not independently checked.} The classical tools marked as
not checked in Section~\ref{sec:tools} (in particular Darroch's original
paper, Graham--Pollak, Weyl and Newton at their sources); the second route
of Remark~\ref{rem:AMW}; the agreement of the Lean statements with the
notes; this manuscript text as a whole.

\section*{Attribution}

\begin{itemize}
\item \emph{Research direction:} Gustavo Bernardo Brito (UFG), project
director: chose the project and the problem, set the objectives and budgets
of the sessions (for this work: to resolve the upper bound). No mathematical
input.
\item \emph{Human mathematical contribution:} none.
\item \emph{AI-generated proof strategies and mathematical discoveries:}
\nolinkurl{claude-opus-5-5} (Lemmas S40 to S47, Theorems T9 and T10,
Corollary S48), 27 November 2026.
\item \emph{Computational contribution:} \nolinkurl{claude-opus-5-5} wrote and
ran \nolinkurl{upper_bound.py} (Python 3.12.4, python-flint 0.9.0, networkx 3.7;
the hardware is not recorded in the \texttt{run} block of the result file).
\item \emph{Formalization:} \nolinkurl{claude-opus-5-5} (formalization agent).
\item \emph{Independent verification:} the four AI reviews listed above, run
as fresh subagent contexts; none is human review.
\item \emph{External mathematical results:} \cite{GL78,GP71,Dar64,Pit97,Mer90,Col89,AAB,AMW}
and the textbook facts of Section~\ref{sec:tools}.
\end{itemize}
Authorship of an eventual manuscript is to be decided by the director under
the project's attribution policy; operating an AI system does not by itself
confer authorship, and AI systems are not listed as authors.

\section*{Limitations}

\begin{itemize}
\item The proofs are candidate proofs checked by AI reviews only; no human
has read them. AI review is not peer review.
\item The proof uses classical results that were cited, not reproved, and
some were not checked at their sources (Section~\ref{sec:tools}).
Section~\ref{sec:path} depends on Collins's formula (C), checked against the
distance matrix only for $n\le30$.
\item Theorem~\ref{thm:T10} is very probably known; the novelty of the other
results is not established, and the literature search has the gaps listed
in Section~\ref{sec:literature}.
\item Only one sequence lemma is formalized, under hypotheses; its receipt
was produced on a working tree that was not clean.
\item The lower half of the conjecture is not treated here;
Corollary~\ref{cor:both} is conditional on it.
\item This \LaTeX{} text was assembled by an AI system from the project notes
and has not itself been reviewed; discrepancies with the notes, if any, are
errors of this draft.
\end{itemize}

\begin{thebibliography}{99}

\bibitem{AAB}
G.~Aalipour, A.~Abiad, Z.~Berikkyzy, L.~Hogben, F.~H.~J. Kenter, J.~C.-H. Lin
and M.~Tait, \emph{Proof of a conjecture of Graham and Lov\'asz concerning
unimodality of coefficients of the distance characteristic polynomial of a
tree}, arXiv:1507.02341. Journal version reported as Electron.\ J.\ Linear
Algebra 34 (2018) 373--380 (not checked against the journal).

\bibitem{AMW}
E.~O.~D. Andriantiana, V.~Razanajatovo Misanantenaina and S.~Wagner,
\emph{The average size of matchings in graphs}, Graphs Combin.\ 36 (2020)
539--560, doi:10.1007/s00373-020-02136-1; arXiv:1903.09825 (version~1 read).

\bibitem{BP}
A.~Berman and R.~J. Plemmons, \emph{Nonnegative Matrices in the
Mathematical Sciences}, SIAM, Philadelphia, 1994.

\bibitem{BH}
A.~E. Brouwer and W.~H. Haemers, \emph{Spectra of Graphs}, Springer, New
York, 2012.

\bibitem{Col89}
K.~L. Collins, \emph{On a conjecture of Graham and Lov\'asz about distance
matrices}, Discrete Appl.\ Math.\ 25 (1989) 27--35,
doi:10.1016/0166-218X(89)90044-9.

\bibitem{Dar64}
J.~N. Darroch, \emph{On the distribution of the number of successes in
independent trials}, Ann.\ Math.\ Statist.\ 35 (1964) 1317--1321.

\bibitem{GL78}
R.~L. Graham and L.~Lov\'asz, \emph{Distance matrix polynomials of trees},
Adv.\ Math.\ 29 (1978) 60--88.

\bibitem{GP71}
R.~L. Graham and H.~O. Pollak, \emph{On the addressing problem for loop
switching}, Bell System Tech.\ J.\ 50 (1971) 2495--2519.

\bibitem{HLP}
G.~H. Hardy, J.~E. Littlewood and G.~P\'olya, \emph{Inequalities}, 2nd ed.,
Cambridge University Press, 1952.

\bibitem{HJ}
R.~A. Horn and C.~R. Johnson, \emph{Matrix Analysis}, 2nd ed., Cambridge
University Press, 2013.

\bibitem{Mer90}
R.~Merris, \emph{The distance spectrum of a tree}, J.\ Graph Theory 14
(1990) 365--369.

\bibitem{Pit97}
J.~Pitman, \emph{Probabilistic bounds on the coefficients of polynomials
with only real zeros}, J.\ Combin.\ Theory Ser.\ A 77 (1997) 279--303.

\end{thebibliography}

\end{document}
